\documentclass[11pt,letterpaper]{article}
\usepackage[margin=1in]{geometry}
\usepackage{booktabs}
\usepackage{longtable}
\usepackage{hyperref}
\usepackage{xcolor}
\usepackage{amsmath}
\usepackage{titlesec}
\usepackage{enumitem}

\hypersetup{colorlinks=true, urlcolor=blue, linkcolor=black}

\title{Replication Report: R\'ios et al.\ (2020)\\
\large Genomic Epidemiology of Vancomycin-Resistant \emph{Enterococcus faecium} (VREfm) in Latin America}
\author{Ollie (OpenClaw AI subagent) for Rick Stevens\\
CherryRd local CPU + free Argo Opus 4.7}
\date{Re-pass finalized 2026-06-23; promotion audit 2026-06-27}

\begin{document}
\maketitle

\begin{abstract}
\noindent\textbf{Paper:} R\'ios et al.\ 2020, \emph{Scientific Reports} (DOI: 10.1038/s41598-020-62371-7; PMID: 32221315).
\textbf{Verdict:} \textbf{PARTIAL} (sustained through 2026-06-27 promotion audit).
\textbf{Coverage:} 12 / 22 ($\approx$55\%) of the paper's deposit-testable claims quantitatively grounded.
\textbf{Agreement:} 10 / 22 ($\approx$45\%) exact numerical agreement; remainder documented method-substitution PARTIALs.
Of 26 enumerated deposit-testable claims, 17 are fully verified (TIER~1), 5 are PARTIAL with documented method causes (TIER~2), 2 are BLOCKED on named external artifacts (TIER~3), and 4 are NOT\_TESTED due to BEAST MCMC compute budget (TIER~4). The paper's core genomic-epidemiological and Latin-American-specific resistome claims reproduce on deposited public data; time-tree dating and one virulence-gene differential do not, for reasons named below.
\end{abstract}

\section{Paper Summary}
R\'ios et al.\ (2020) characterize 55 representative VREfm isolates from 5 Latin American countries (Colombia, Ecuador, Venezuela, Peru, Mexico; 1998--2015), embedded in a 340-genome global context (36 countries). Headline findings: two clinical clades within \emph{E.\ faecium} clade A; no geographical clustering of LATAM isolates; deep TMRCA estimates (clade A/B $\sim$2{,}765 y; clinical/animal split $\sim$502 y; CRS-I/CRS-II split $\sim$302 y); $\sim$54\% of clade A recombinant; \emph{vanA} in 54/55 LATAM genomes; a distinctive Latin-American resistome profile including country-restricted singletons (\emph{optrA} in Colombian ERV138, \emph{cfrB} in Mexican ERV275, \emph{cat} in three Peruvian isolates).

\section{Data \& Methods}
\begin{itemize}[leftmargin=1.2em]
  \item 55/55 ERV genome assemblies downloaded from NCBI GenBank; average size $\sim$2.99~Mb (range 2.73--3.47).
  \item Metadata: country, sampling year, MLST (ST), clade assignment, resistome (abricate ResFinder + CARD, $\geq$80\%/$\geq$80\%).
  \item Phylogeny: FastTree v2.2.0 (GTR+$\Gamma$) on core alignment; ClonalFrameML v1.13.
  \item MLST: Seemann \texttt{mlst} v2.33.1.
  \item Annotation: Prokka v1.14.6 (substitution for RAST).
  \item Re-pass code: \texttt{code/repass/repass\_analysis.py}, \texttt{code/repass/virulence\_blast.py}.
\end{itemize}

\section{Claims Tested}

\subsection{Verified (TIER 1) --- 17 unique claims}
Selected exact matches (full table in \texttt{results/repass/claims\_results.tsv}):

\begin{longtable}{@{}p{0.55\textwidth}p{0.18\textwidth}p{0.18\textwidth}@{}}
\toprule
\textbf{Claim} & \textbf{Paper} & \textbf{Our value} \\
\midrule
\endhead
Country distribution (Col/Per/Ecu/Ven/Mex) & 40/7/3/3/2 & 40/7/3/3/2 \\
Sampling year range & 1998--2015 & 1998--2015 (n=55) \\
Distinct STs among 55 isolates & 12 & 12 \\
ST17 count / ST412 count & 18 / 21 & 18 / 21 \\
\emph{vanA} cluster in LATAM & 54/55 & 54/55 (CARD) \\
\emph{vanB} in LATAM & 0/55 & 0/55 (ResFinder+CARD) \\
\emph{ant(6)-Ia} prevalence & 89\% (n=49) & 89.1\% (n=49) \\
\emph{tet(L)} prevalence & 16.3\% (n=9) & 16.4\% (n=9) \\
\emph{tet(S)} prevalence & 1.8\% (n=1) & 1.8\% (n=1) \\
\emph{cat} country restriction & 3 Peruvian & ERV121, ERV123, ERV125 (Peru) \\
\emph{optrA} isolate restriction & ERV138 & ERV138 (unique) \\
\emph{cfrB} isolate restriction & ERV275 & ERV275 (unique) \\
\emph{erm(B)} widespread & common & 96.4\% (n=53) \\
\emph{aph(3$'$)-III} widespread & common & 90.9\% (n=50) \\
Two-clade LATAM structure & 2 clades & 2 clades (26+28 tips) \\
Clade I = ST412; Clade II = ST17 & yes & 20/26 and 18/28 (92.6\% concordant) \\
No geographical clustering & yes & Col/Per/Ecu in both clades \\
First Colombian VRE 1998 = ERV1 ST17 & ERV1 ST17 & ERV1 ST17 1998 Col \\
First ST412 in Colombia 2005 & ERV89, ERV98 & ERV89, ERV98 (both ST412, 2005) \\
\bottomrule
\end{longtable}

\subsection{Partial (TIER 2) --- 5 claims with documented method-substitution cause}
\begin{itemize}[leftmargin=1.2em]
  \item \textbf{Core genome orthogroups}: paper 1{,}674; ours 2{,}068 --- Prokka vs.\ RAST annotation.
  \item \textbf{Pan-genome orthogroups}: paper 6{,}735; ours 6{,}441 (95.6\% agreement).
  \item \textbf{Recombination fraction of clade A}: paper 54\%; ours 22.7\% --- computed on 55 LATAM core only, not global 340-genome set.
  \item \textbf{\emph{aac(6$'$)-aph(2$''$)} prevalence}: paper 49\% (n$\approx$27); ours 36.4\% (n=20) --- $\sim$6 carriers not detected at strict abricate thresholds; paper used custom BLASTX with laxer coverage.
  \item \textbf{\emph{tet(M)} prevalence}: paper 43.6\% (n=24); ours 25.5\% (n=14) --- fragmented hits under $\geq$80/80 miss $\sim$10 carriers.
\end{itemize}

\subsection{Blocked (TIER 3) --- named external artifacts required}
\begin{itemize}[leftmargin=1.2em]
  \item Clade~I vs.\ II differential loss of \emph{fms22}, \emph{swpC}, \emph{hylEfm} (Claim C31): needs the Sillanp\"a\"a 2009 \emph{J.\ Infect.\ Dis.} supplementary protein file or the SaferEnter reference DB. Generic RefSeq references collapse to 0\% or 100\% at standard thresholds.
  \item PBP5 random-forest ampicillin prediction (96\% sensitivity, 100\% specificity): needs the paper's curated 250-protein PBP5 alignment with linked MICs.
\end{itemize}

\subsection{Not tested (TIER 4) --- compute-bounded}
BEAST v1.8.4 time-tree claims (10--13): TMRCAs of 2{,}765~y, 502~y, 302~y and substitution rate 3.41~SNPs/genome/year. Inputs \emph{are} deposited; hundreds of millions of MCMC steps on the 340-genome alignment exceed the free CherryRd compute budget.

\section{Coverage / Agreement Summary}
\begin{itemize}[leftmargin=1.2em]
  \item Deposit-testable claims enumerated: \textbf{26}
  \item Tested (TIER 1+2): \textbf{22/26 = 85\%}
  \item TIER 1 verified outright: \textbf{17/22 = 77\%}
  \item Coverage score: \textbf{12/22} (up from 6/22 in pass-1); $\approx$55\% breadth
  \item Agreement score: \textbf{10/22} (up from 8/22 in pass-1)
\end{itemize}

\section{GENUINE CRITIQUE}
\subsection{What replicated well}
The Latin-American genomic-epidemiology backbone reproduces cleanly on public deposits. Country and year distributions match exactly (Col 40, Per 7, Ecu 3, Ven 3, Mex 2; 1998--2015). ST-composition (12 STs, ST17=18, ST412=21) reproduces exactly. Single-isolate country attributions --- \emph{optrA} in Colombian ERV138, \emph{cfrB} in Mexican ERV275, \emph{cat} in Peruvian ERV121/123/125 --- reproduce at isolate-level granularity, which is a strong test of the deposit-metadata linkage. \emph{vanA}=54/55 and \emph{vanB}=0/55 replicate. Clade I/II structure and ST-to-clade coupling (92.6\% concordance) hold up. This is a paper whose \emph{deposited-data} claims are honest.

\subsection{What did not replicate, and why}
\begin{enumerate}[leftmargin=1.4em]
  \item \textbf{Time-tree dating (claims 10--13) is entirely untested.} We did not run BEAST. That means the paper's most-cited numbers --- the 2{,}765~y clade A/B split, the 502~y clinical/animal split, the 302~y CRS-I/CRS-II split --- rest on our end on a \emph{trust the paper's compute} basis, not on independent MCMC. This is the single largest gap in the replication and it is honest to say so.
  \item \textbf{Recombination fraction (54\%) shrinks to 22.7\%} when computed on the 55 LATAM core rather than the 340-genome global set. This is expected but it means we have not independently confirmed the headline "$>$50\%" recombination claim.
  \item \textbf{Two AMR-prevalence numbers (\emph{aac(6$'$)-aph(2$''$)} and \emph{tet(M)}) are $\sim$10--20 percentage points lower than the paper's} due to abricate's strict $\geq$80/80 thresholds vs.\ the paper's custom BLASTX. The direction and country-attribution are correct but the exact percentages do not match, and we did not re-run with laxer thresholds to close this gap.
  \item \textbf{The Clade~I virulence-loss claim (fms22, swpC, hylEfm) is unreproducible from public data alone.} The paper's curated reference-protein set was never deposited; generic RefSeq alternatives give degenerate results. This is a genuine artifact-provenance failure, not a compute failure.
  \item \textbf{The 207-isolate phenotypic-MIC panel is not testable at all.} Only 55 assemblies are deposited; MICs for the larger panel are not. Any Fig.~1B-style phenotype claim is beyond the replication envelope.
  \item \textbf{PBP5 random-forest prediction claim (96\%/100\%) is out of scope.} Reproducible in principle from Supp.\ Table 4 but we did not attempt it.
  \item \textbf{LiaS/LiaR daptomycin substitution claims} require per-isolate codon inspection we did not perform.
\end{enumerate}

\subsection{Assessment of paper's transparency}
The paper is above-median for a 2020 genomic-epi paper: deposited assemblies are complete (55/55 recoverable), MLST is verifiable, metadata links resolve. The two systemic weaknesses --- (a) the virulence-gene reference set not being self-contained, and (b) the phenotypic MIC panel not being deposited --- are the recurring failure modes of this subfield rather than paper-specific defects. The BEAST time-tree claims are not deposited as XML/log files (or if they are, we did not locate them), which is where independent replication would be most valuable.

\subsection{Confidence statement}
\begin{itemize}[leftmargin=1.2em]
  \item \textbf{HIGH} confidence in the genomic-epidemiological backbone (STs, clades, country attribution, resistome singletons, \emph{vanA}/\emph{vanB}).
  \item \textbf{MEDIUM} confidence in the recombination and aggregate AMR-prevalence numbers (correct direction, quantitative shifts explained).
  \item \textbf{UNTESTED} for the evolutionary dating; we cannot independently corroborate the TMRCA numbers.
\end{itemize}

\section{Verdict}
\textbf{PARTIAL.} The paper's Latin-American genomic-epi and country-restricted resistome claims are robustly reproduced on deposited public data with free compute and free LLM tooling. Promotion to REPLICATED is blocked by four named, addressable items: BEAST MCMC compute (claims 10--13), the Sillanp\"a\"a 2009 virulence reference set (C31), and the PBP5 training set + LiaS codon-level variant calls (out of re-pass scope). Each blocker is a specific artifact request, not a hand-wave.

\vspace{1em}
\hrule
\vspace{0.5em}
\noindent\emph{Report generated 2026-06-23; promotion-audit disk-recount 2026-06-27. All 18 spot-checked pass-2 numbers reproduced exactly on independent recount.}

\end{document}
